\chapter*{The Argument of Part III}
\pdfbookmark[0]{The Argument of Part III}{part3.guide}
\label{part3:guide}
\markboth{Open Source Evaluatology}{The Argument of Part III}
\noindent Parts I and II established what an evaluation must specify and how collaboration can be measured. The present part asks what those commitments require in use. Its progression follows three linked objects: a contribution considered for recognition, a capability whose continuity matters, and an arrangement that distributes work, authority, benefits, and burdens across an ecosystem.

The chapters share a conceptual framework, but each retains its own research design. The course studies examine evidence and judgment within educational assessment. The archive study examines reasons, trajectories, and practices around discontinuation. The Redis study follows participation and organizational relationships around a governance change. Their findings accumulate through the distinctions they establish; their samples and outcomes are not pooled.

\vspace{4pt}
\begingroup\small
\renewcommand{\arraystretch}{1.16}
\noindent\begin{tabularx}{\linewidth}{@{}>{\raggedright\arraybackslash}p{30mm}>{\raggedright\arraybackslash}X >{\raggedright\arraybackslash}X@{}}
\toprule
\textbf{Object} & \textbf{Principal evidence} & \textbf{Refinement of the framework}\\
\midrule
\textbf{7\quad Contribution} & Two course studies: a published 2023 comparison and an unpublished 2025 assessment deployment. & Observation depends on opportunity; relational evidence guides inspection; correction can change the basis of a decision.\\[5pt]
\textbf{8\quad Continuity} & The doctoral study of 361 archived software projects and its question-specific subsets. & Follow the function, responsible parties, and transition dates beyond the original repository.\\[5pt]
\textbf{9\quad Governance} & Redis and Valkey, selected licensing-change comparisons, and a developer-distribution analysis. & Distinguish contribution, affiliation, authority, and affected benefit; preserve the rules of attribution and aggregation.\\
\bottomrule
\end{tabularx}
\endgroup
\par\vspace{9pt}
Chapter 7 connects recognition to the review, maintenance, and coordination work that Chapter 8 follows over time. Chapter 8 connects a possible successor to the responsibilities and functions it must sustain. Chapter 9 then examines participation across a fork while adding organizational relationships and decision rights. A change of analytical boundary creates new evidential requirements at each step.

The measurement language remains that of Part II. \emph{Activity} names a defined measure; event counts and commit trajectories retain their own definitions. Global influence, local contribution, and application-specific OpenRank configurations remain distinct. Historical findings retain the versions used to produce them. A later formulation can motivate a new analysis without changing what an earlier result meant.

The return to theory is equally specific. A contribution record is conditional on its observation setting; continuity belongs to a function under stated conditions; and governance judgment requires criteria as well as participation evidence. Part IV carries these requirements into think-tank applications and an open tool ecosystem, then examines how the resulting claims can be validated, compared, and incorporated into a cumulative discipline. The studies here establish the questions that this institutional and methodological development must answer.
\clearpage
